\section{Appendix: An example of database usage}

\textbf{MOVE THIS TO ONLINE-ONLY?}

Here I will work through an example of a real scientific question using several database systems. I will focus on NoSQL databases, for two reasons:

\begin{enumerate}
\item They are less well known amongst scientists compared to relational databases
\item I personally prefer the NoSQL approach in most cases due to their flexibility
\end{enumerate}

Note that while many researchers may never use all the different databases used here, I am intentionally presenting a full-stack example to show how they can be combined.

The question that I will ask is as follows: How well can the biological similarity between traits (including diseases as well as other phenotypes) be estimated from the semantic similarity of publications that refer to the trait?  I will use two primary datasets to assess this:

\begin{itemize}
\item A dataset of genome-wise association study (GWAS) results for specific traits obtained from \href{https://www.ebi.ac.uk/gwas/docs/file-downloads}{here}.
\item Abstracts that refer to each of the traits identified in the GWAS result, obtained from the \href{https://pubmed.ncbi.nlm.nih.gov/}{PubMed} database.
\end{itemize}

I will not present all of the code for each step, which can be found in the project reposistory (github.com/BetterCodeBetterScience/example-databases). Rather, I will show portions that are particularly relevant to the databases being used.

\subsection{Adding GWAS data to a document store}

We start by uploading the GWAS data and adding them to a document store, using \textit{MongoDB}, which I installed on my local machine.   Upon reading the CSV file containing the data and examining their structure, we see that they are not properly \textit{normalized}. Normalization is a concept that is essential for relational databases but can also be very helpful for document stores.  We won't go into the details of normalization here; rather, we will simply outline the primary requirements of the \textit{first normal form}, which is the most basic form of normalization.  This requires that:\index{MongoDB}

\begin{enumerate}
\item There are no duplicate records
\item The columns contain scalar values (and thus do not contain composite values such as sets of values)
\item There are not multiple columns that contain the same kind of data
\end{enumerate}

In this case, looking at the data we see that several columns contain multiple values for genes, due to the fact that some single nucleotide polymorphisms (SNPs) are located in intergenic space and thus can be mapped to multiple genes.  These values separated by commas (e.g., ``FADS2, FADS1'').  To normalize this, we can \textit{explode} the data frame, which involves separating out these values into separate rows, which have otherwise identical contents.

\begin{lstlisting}[style=Pythonshort]
gwas_data = get_exploded_gwas_data()
\end{lstlisting}

We can now import the data from this data frame into a MongoDB collection, mapping each unique trait to the genes that are reported as being associated with it.  First I generated a separate function that sets up a MongoDB collection.  We can then use that function to set up our gene set collection:

\begin{lstlisting}[style=Python]
# since the collection names are shared across functions, 
# create global variable with their names
COLLECTION_GENESET = 'geneset_annotations_by_trait'

def import_genesets_by_trait(
    gwas_data_melted: pd.DataFrame
) -> None:

    geneset_annotations_by_trait = setup_mongo_collection(COLLECTION_GENESET)
    # "mapped_trait_uri" is the field that contains the identifier for each trait
    geneset_annotations_by_trait.create_index('mapped_trait_uri', unique=True)

    # first get a mapping from MAPPED_TRAIT_URI to TRAIT_NAME
    trait_name_mapping = gwas_data_melted.set_index('MAPPED_TRAIT_URI')[
        'MAPPED_TRAIT'
    ].to_dict()

    # loop through each unique MAPPED_TRAIT_URI in gwas_data data frame 
    # add all of its gene sets to the mongo collection 
    for mapped_trait_uri in tqdm(
        gwas_data_melted['MAPPED_TRAIT_URI'].unique()
    ):
        gene_sets = (
            gwas_data_melted[
                gwas_data_melted['MAPPED_TRAIT_URI'] == mapped_trait_uri
            ]['GENE_ID']
            .unique()
            .tolist()
        )
        # clean up gene names
        gene_sets = [gene.strip() for gene in gene_sets]
        # add the item to the mongodb collection, also
        # including the trait name
        geneset_annotations_by_trait.update_one(
            {'mapped_trait_uri': str(mapped_trait_uri)},
            {
                '$set': {
                    'mapped_trait_uri': str(mapped_trait_uri),
                    'gene_sets': gene_sets,
                    'trait_name': trait_name_mapping.get(mapped_trait_uri, ''),
                }
            },
            upsert=True,
        )


import_genesets_by_trait(gwas_data)
\end{lstlisting}

Note that this function uses an \textit{upsert} operation, which is a combination of insertion (if the document with this key doesn't already exist) or updating (if the document with this key already exists).  This results in a collection with 3,047 records, each of which is indexed by the trait identifier and includes a list of all of the genes associated with the trait across the GWAS studies.  We also include the trait names to make the records human-readable, but we rely upon the unique identifiers which in this case are defined as URLs, such as ``http://www.ebi.ac.uk/efo/EFO\_0004309'' which maps to the trait of ``platelet count''.

\subsection{Annotating gene sets}

Remember that our goal in this analysis is to identify the biological overlap between traits. We could do this by assessing the degree to which they are associated with the same genes, but this would miss out on the fact that genes work together in networks that are often associated with the function of specific biological pathways.  Given a particular set of genes, we can use bioinformatics tools to identify the biological processes that are associated with that set of genes. In this case I used the  \textit{g:Profiler} tool from ELIXIR, which comes with a handy Python package.   This tool returns a set of pathways that are statistically enriched for each gene set, each of which is defined by a unique identifier that refers to a particular ontology such as the Gene Ontology.   For example, the 877 genes associated with the ``platelet count'' trait are involved in a range of biological processes and pathways, which range from the very general (e.g., GO:0005515, referring to ``protein binding'') to the very specific (e.g., GO:0007599, referring to ``hemostasis'', which is the stopping of bleeding).

We can use \textit{g:Profiler} to obtain the annotation for the gene set associated with each trait:

\begin{lstlisting}[style=Python]
def annotate_genesets_by_trait() -> None:
    # loop over all entries in the geneset_annotations_by_trait collection
    # and do functional annotation of the gene sets

    geneset_annotations_by_trait = setup_mongo_collection(COLLECTION_GENESET)

    gp = GProfiler(return_dataframe=True)

    # use a list here so that we can use tqdm to show progress
    # skip any entries that already have functional_annotation
    annotations = [
        i
        for i in geneset_annotations_by_trait.find({})
        if 'functional_annotation' not in i
    ]

    for entry in tqdm(annotations):
        mapped_trait_uri = entry['mapped_trait_uri']
        gene_sets = entry['gene_sets']
        if len(gene_sets) == 0:
            continue
        # do functional annotation
        try:
            annotation_results = gp.profile(
                organism='hsapiens',
                query=gene_sets,
                sources=['GO:BP', 'GO:MF', 'GO:CC', 'KEGG', 'REAC'],
            )
        except Exception as e:
            # catch any exception to avoid breaking the loop
            print(f'Error annotating {mapped_trait_uri}: {e}')
            continue

        # convert the dataframe to a dictionary
        annotation_results_dict = annotation_results.to_dict(orient='records')
        # update the entry in the mongo collection with the annotation results
        geneset_annotations_by_trait.update_one(
            {'mapped_trait_uri': str(mapped_trait_uri)},
            {'$set': {'functional_annotation': annotation_results_dict}},
        )
    # drop members of geneset_annotations_by_trait with empty annotation
    geneset_annotations_by_trait.delete_many(
        {'functional_annotation': {'$in': [None, [], {}]}}
    )

    print(
        f'Remaining entries with functional annotation: {
          geneset_annotations_by_trait.count_documents({})}'
    )


annotate_genesets_by_trait()
\end{lstlisting}

\begin{lstlisting}[style=shell, numbers=none]
Number of documents in geneset_annotations_by_trait: 3047
100%|██████████| 3047/3047 [56:31<00:00,  1.11s/it] 
Remaining entries with functional annotation: 1845
\end{lstlisting}

Having now annotated the gene sets with information about their biological functions, we can now move to assessing the similarity of traits based on their associated functions.

\subsection{Mapping pathway information to traits}

The previous analysis added a \texttt{functional annotation} element to each trait, which includes information about the associated pathways. While we could use this collection to obtain the mappings from traits to pathways, the deeply embedded nature of the annotation data would make the queries somewhat complicated.  Instead we will use that information to generate a collection including all unique pathways, which we will then use to compute biological similarity between traits:

\begin{lstlisting}[style=Python]
COLLECTION_PATHWAYS = 'pathways'

def get_pathway_info_by_trait() -> None:
    geneset_collection = setup_mongo_collection(
        COLLECTION_GENESET, clear_existing=False
    )
    pathway_collection = setup_mongo_collection(
        COLLECTION_PATHWAYS, clear_existing=False
    )

    # loop through traits and add pathway information to the database
    traits = [
        i
        for i in geneset_collection.find()
        if 'functional_annotation' in i and len(i['functional_annotation']) > 0
    ]

    for trait in tqdm(traits):
        annotations = trait['functional_annotation']
        for pathway in annotations:
            # change key for clarity
            pathway['pathway_id'] = pathway.pop('native')

            pathway_collection.update_one(
                {'pathway_id': pathway['pathway_id']},
                {
                    '$set': {
                        'name': pathway.get('name', ''),
                        'source': pathway.get('source', ''),
                        'description': pathway.get('description', ''),
                    }
                },
                upsert=True,
            )


get_pathway_info_by_trait()
\end{lstlisting}

\begin{lstlisting}[style=shell, numbers=none]
Number of documents in geneset_annotations_by_trait: 1845
Number of documents in pathways: 0
100%|██████████| 1845/1845 [00:18<00:00, 98.22it/s]
\end{lstlisting}

\subsection{Generate the graph database linking publications pathways to traits}

Graph databases are designed to store and query relational information, making them perfect for our dataset linking pathways to traits.  In the next step I created a graph database using the Neo4j database package, which involves creating a Neo4j session and then issuing a set of Cyper commands:

\begin{lstlisting}[style=Python]
def get_neo4j_session():
    assert 'NEO4J_PASSWORD' in os.environ, 'NEO4J_PASSWORD should be set in .env'
    neo4j_driver = GraphDatabase.driver(
        'bolt://localhost:7687', auth=('neo4j', os.environ['NEO4J_PASSWORD'])
    )
    return neo4j_driver.session()


def build_neo4j_graph() -> None:
    geneset_collection = setup_mongo_collection(COLLECTION_GENESET)
    pathway_collection = setup_mongo_collection(COLLECTION_PATHWAYS)

    with get_neo4j_session() as session:
        # Clear DB
        session.run('MATCH (n) DETACH DELETE n')
        
        # Create Indexes
        session.run((
          'CREATE CONSTRAINT IF NOT EXISTS FOR (p:Phenotype) '
          'REQUIRE p.id IS UNIQUE'))
        session.run(
          ('CREATE CONSTRAINT IF NOT EXISTS FOR (p:Pathway) '
          'REQUIRE p.id IS UNIQUE'))

        # 1. Batch Import Pathways
        # Fetch all pathways into a list of dicts
        print("Loading pathways...")
        pathways_data = list(pathway_collection.find({},
          {'_id': 0, 'pathway_id': 1, 'name': 1, 'source': 1, 'description': 1}))
        
        # Use UNWIND to insert all pathways in one transaction
        session.run("""
            UNWIND $batch AS row
            MERGE (pw:Pathway {id: row.pathway_id})
            SET pw.name = row.name,
                pw.source = row.source,
                pw.description = row.description
        """, batch=pathways_data)

        # 2. Batch Import Phenotypes and Relationships
        print("Loading phenotypes and relationships...")
        # We need to restructure the Mongo data slightly for Neo4j consumption
        pheno_batch = []
        for doc in geneset_collection.find():
            pheno_id = str(doc['mapped_trait_uri'])
            # Extract list of pathway IDs
            pathway_ids = [
                i['native'] 
                for i in doc.get('functional_annotation', []) 
                if 'native' in i]
             
            if pathway_ids:
                pheno_batch.append({
                    'id': pheno_id,
                    'name': doc.get('trait_name', ''),
                    'pathway_ids': pathway_ids
                })

        # Insert Phenotypes and create edges to Pathways
        session.run("""
            UNWIND $batch AS row
            MERGE (p:Phenotype {id: row.id})
            SET p.name = row.name
            
            WITH p, row
            UNWIND row.pathway_ids AS pw_id
            MATCH (pw:Pathway {id: pw_id})
            MERGE (p) -[:MAPPED_TO] ->(pw)
        """, batch=pheno_batch)

    print("Graph build complete.")


build_neo4j_graph()
\end{lstlisting}

\begin{lstlisting}[style=shell]
Number of documents in geneset_annotations_by_trait: 1845
Number of documents in pathways: 6051
Loading pathways...
Loading phenotypes and relationships...
Graph build complete.
\end{lstlisting}

Having populated the database, we can then use the \texttt{gds.nodeSimilarity} function from the Neo4j Graph Data Science library to compute the similarity of each trait in terms of their pathway overlap.  This function takes a \textit{bipartite graph} (which is a graph that contains two sets of nodes - in our case, pathways and traits) and returns the Jaccard similarity coefficient, which in this case reflects the ratio of common pathways to the total number of pathways for each pair of nodes. We filter out nodes that have fewer than two pathways associated with them to make the estimate more stable. (Note that here I use the term ``phenotype'' interchangeably with ``trait''.)

\begin{lstlisting}[style=Python]
def compute_phenotype_similarities(
    graph_name: str = 'phenotype -pathway -graph',
    min_pathways: int = 2
) -> pd.DataFrame:

    with get_neo4j_session() as session:
        # 1. Clean up any existing graph with the same name
        try:
            session.run("CALL gds.graph.drop($name)", name=graph_name)
        except ClientError:
            # Graph did not exist, safe to ignore
            pass

        # 2. Project the Graph
        # We project Phenotypes, Pathways, and the relationship between them.
        # 'UNDIRECTED' orientation allows the algorithm to see the shared connections.
        session.run("""
            CALL gds.graph.project(
                $graph_name,
                ['Phenotype', 'Pathway'],
                {
                    MAPPED_TO: {
                        orientation: 'UNDIRECTED'
                    }
                }
            )
        """, graph_name=graph_name)

        # 3. Run Node Similarity
        # We use 'degreeCutoff' to filter out phenotypes with too few pathways 
        # BEFORE the calculation - much faster than filtering the results.
        result = session.run("""
            CALL gds.nodeSimilarity.stream($graph_name, {
                degreeCutoff: $min_pathways
            })
            YIELD node1, node2, similarity
            
            // Map internal IDs back to our Phenotype IDs
            WITH node1, node2, similarity
            MATCH (p1:Phenotype), (p2:Phenotype)
            WHERE id(p1) = node1 AND id(p2) = node2
            
            RETURN p1.id AS phenotype1, p2.id AS phenotype2, similarity
            ORDER BY similarity DESC
        """, graph_name=graph_name, min_pathways=min_pathways)

        # 4. Convert to DataFrame
        # result.data() fetches all records into a list of dictionaries
        df = pd.DataFrame(result.data())
        
        # 5. Cleanup: Drop the graph to free up memory
        try:
            session.run("CALL gds.graph.drop($name)", name=graph_name)
        except ClientError:
            pass

        return df


similarity_result_df = compute_phenotype_similarities()
\end{lstlisting}

Now that we have the similarity between phenotypes based on their biological functions, we can move on to assessing their similarity based on associated publications.

\subsection{Obtaining literature related to traits}

To assess semantic similarity between traits we need to obtain abstracts related to each trait.  To do this, we first wish to obtain any synonyms for the traits, to maximize the effectiveness of the search.  We can obtain these from the EBI Ontology Search API (www.ebi.ac.uk/ols4/api-docs):

\begin{lstlisting}[style=Python]
COLLECTION_TRAIT_INFO = 'trait_info_by_trait'

def get_trait_info_from_ols(
    client_url: str = 'http://www.ebi.ac.uk/ols',
) -> None:
    # use EBI OLS API to get trait information for all traits
    trait_info_by_trait = setup_mongo_collection(
        collection_name=COLLECTION_TRAIT_INFO
    )

    trait_info_by_trait.create_index('trait_uri', unique=True)

    geneset_annotations_by_trait = setup_mongo_collection(
        collection_name=COLLECTION_GENESET
    )

    # get all unique trait URIs that are not already in the collection
    # use lstrip because many have a leading space
    unique_trait_uris = [
        i.lstrip()
        for i in geneset_annotations_by_trait.distinct('mapped_trait_uri')
        if trait_info_by_trait.count_documents({'trait_uri': i.lstrip()}) == 0
    ]
    print(f'Found {len(unique_trait_uris)} un -annotated trait URIs.')

    client = ols_client.Client(client_url)

    for trait_uri in tqdm(unique_trait_uris):
        trait_id = trait_uri.split('/')[ -1]
        trait_uri = str(trait_uri)
        # skip if already in the collection
        if trait_info_by_trait.count_documents({'trait_uri': trait_uri}) > 0:
            continue
        try:
            term_data = get_info_from_ols(trait_id, client)
        except HTTPError:
            print(f'HTTPError for {trait_uri}')
            continue
        if term_data is None:
            print(f'No data returned for {trait_uri}')
            continue
        trait_info_by_trait.update_one(
            {'trait_uri': str(trait_uri)},
            {'$set': {'trait_uri': str(trait_uri), 'trait_info': term_data}},
            upsert=True,
        )


get_trait_info_from_ols()
\end{lstlisting}

\begin{lstlisting}[style=shell, numbers=none]
Number of documents in geneset_annotations_by_trait: 1845
Number of documents in pathways: 0
100%|██████████| 1845/1845 [00:18<00:00, 98.22it/s]
\end{lstlisting}

We then use the \textit{Biopython.Entrez} module to search PubMed for each of the traits, along with any synonyms, obtaining a maximum of 100 abstracts for each query:\index{Biopython}

\begin{lstlisting}[style=Python]
COLLECTION_PMID_BY_TRAIT = 'pmids_by_trait'

def get_pmids_for_traits(
    n_abstracts_per_trait: int = 100, verbose: bool = False
) -> None:

    pmid_collection = setup_mongo_collection(COLLECTION_PMID_BY_TRAIT)
    _ = pmid_collection.create_index(
        [('trait_uri', pymongo.ASCENDING)], unique=True
    )

    # get all entries from the trait_info_by_trait collection and pull out 
    # the label and synonyms to use as pubmed search terms
    trait_info_collection = setup_mongo_collection(COLLECTION_TRAIT_INFO)
    db_result = list(trait_info_collection.find({}))

    for result in tqdm(db_result, desc='Searching PubMed'):
        trait_uri = result['trait_uri']
        lbl = result['trait_info']['label']
        synonyms = result['trait_info'].get('synonyms', [])
        # create a pubmed query using the label and synonyms
        query_terms = [lbl] + synonyms
        query = ' OR '.join([f'"{term}"' for term in query_terms])

        existing_entry = pmid_collection.find_one({'trait_uri': trait_uri})
        # Skip existing entry if pmid entry is not empty
        if (
            existing_entry is not None
            and existing_entry.get('pmids')
            and len(existing_entry.get('pmids')) > 0
        ):
            if verbose:
                print(f'PMIDs already exist for {lbl}, skipping...')
            continue
        # run pubmed search - retry up to 3 times if it fails
        pmids = run_pubmed_search(query, retmax=n_abstracts_per_trait,
          max_retries=3)
        if pmids is None or len(pmids) == 0:
          continue
        pmid_collection.update_one(
            {'trait_uri': trait_uri},
            {'$set': {'label': lbl, 'pmids': pmids, 'search_query': query}},
            upsert=True,
        )


get_pmids_for_traits()
\end{lstlisting}

\begin{lstlisting}[style=shell, numbers=none]
Number of documents in pmids_by_trait: 0
Number of documents in trait_info_by_trait: 1590
Searching PubMed: 100%|██████████| 1590/1590 [13:58<00:00,  1.90it/s]
\end{lstlisting}

We then identify all of the unique PubMed IDs (PMIDs) across these queries and fetch the full PubMed record (including the title and abstract, which we will use in our semantic analysis) for each of them, placing them in a new collection:

\begin{lstlisting}[style=Python]
COLLECTION_PUBMED = 'pubmed_abstracts'


def fetch_and_store_pubmed_abstracts(
    batch_size: int = 500
) -> None:

    pubmed_collection = setup_mongo_collection(
        COLLECTION_PUBMED, clear_existing=False
    )
    pubmed_collection.create_index([('PMID', pymongo.ASCENDING)], unique=True)

    # remove any PMIDs already in the pubmed_collection
    existing_pmids = set()
    for entry in pubmed_collection.find({}, {'PMID': 1}):
        existing_pmids.add(entry['PMID'])

    pmids_to_fetch = [
        pmid
        for pmid in get_unique_pmids_from_trait_collection()
        if pmid not in existing_pmids
    ]

    if len(pmids_to_fetch) == 0:
        print('All PMIDs are already fetched. Skipping.')
    else:
        print(f'Fetching {len(pmids_to_fetch)} PMIDs...')

    for i in tqdm(
        range(0, len(pmids_to_fetch), batch_size),
        desc='Fetching PubMed abstracts',
    ):
        batch = pmids_to_fetch[i : i + batch_size]
        pubmed_records = fetch_pubmed_records(batch, retmax=batch_size)
        parsed_records = parse_pubmed_query_result(pubmed_records)
        if not parsed_records:
            print(f'No new records to insert for batch {i // batch_size + 1}.')
            continue
        pubmed_collection.insert_many(parsed_records.values())
        # print(f"Inserted {len(parsed_records)} abstracts")


fetch_and_store_pubmed_abstracts()
\end{lstlisting}

\begin{lstlisting}[style=shell, numbers=none]
Number of documents in pubmed_abstracts: 0
Number of documents in pmids_by_trait: 1590
Fetching 106387 PMIDs...
\end{lstlisting}

\subsection{Add documents to vector database}

We want to use the documents downloaded from PubMed for each trait to compute the semantic similarity between traits.  This is a good application for a \textit{vector database}, which can ingest documents, embed them into a vector space, which we can then use to perform similarity computations between documents based on their vector embeddings.  We will use \textit{ChromaDB} which is a popular open-source vector database.  By default \textit{ChromaDB} uses the \textit{all-MiniLM-L6-v2} embedding model from the \textit{Sentence Transformers} package, but we will instead use the more powerful \textit{text-embedding-3-large} via the OpenAI API.\index{ChromaDB}\index{OpenAI}

\begin{lstlisting}[style=Python]
def get_chromadb_collection(collection_name='pubmed_docs', 
    path='../../data/chroma_data',
    embedding: str = "text -embedding -3 -large"):

    assert 'OPENAI' in os.environ, 'OPENAI API key should be set in .env'
    if embedding is not None:
        embedding_function = embedding_functions.OpenAIEmbeddingFunction(
                    api_key=os.getenv('OPENAI'),
                    model_name=embedding
            )
    else:
        embedding_function = embedding_functions.SentenceTransformerEmbeddingFunction(
            model_name="all -MiniLM -L6 -v2")
    client = PersistentClient(path=path)
    # check if collection exists, if not create it
    if collection_name in [col.name for col in client.list_collections()]:
        return client.get_collection(name=collection_name)
    else:
        print(f'Created new collection: {collection_name}')
        return client.create_collection(name=collection_name, embedding_function=embedding_function)


def add_pubmed_abstracts_to_chromadb(batch_size: int = 5000) -> None:

    pubmed_collection = setup_mongo_collection(
        COLLECTION_PUBMED, clear_existing=False
    )

    collection = get_chromadb_collection()
    # get ids (pmid) and documents (title + abstract) from pubmed_collection
    ids = []
    documents = []
    for entry in pubmed_collection.find({}):
        full_text = entry.get('title', '') + ' ' + entry.get('abstract', '')
        documents.append(full_text)
        ids.append(str(entry['PMID']))

    # exclude ids that are already in the chromadb collection
    existing_ids = set(collection.get(include=[])['ids'])
    ids_to_add = []
    documents_to_add = []
    for i, id_ in enumerate(ids):
        if id_ not in existing_ids:
            ids_to_add.append(id_)
            documents_to_add.append(documents[i])

    # add in batches
    for i in range(0, len(ids_to_add), batch_size):
        batch_ids = ids_to_add[i : i + batch_size]
        batch_documents = documents_to_add[i : i + batch_size]
        collection.add(ids=batch_ids, documents=batch_documents)
        print(f'Added {len(batch_ids)} documents to chromadb collection')


add_pubmed_abstracts_to_chromadb()
\end{lstlisting}

\begin{lstlisting}[style=shell, numbers=none]
Number of documents in pubmed_abstracts: 106120
Using existing collection: pubmed_docs
\end{lstlisting}

We then compute the vector similarity between the document embeddings associated with each pair of traits. We compute the Euclidean distance between each pair of documents associated with both traits, and then compute the mean across the documents.  These are added to the data frame that also contains the pathway distances, for each in subsequent analyses:

\begin{lstlisting}[style=Pythonshort]
text_similarity_df = compute_text_similarities(similarity_result_df)
\end{lstlisting}

We now have the semantic and biological similarity values for each pair of traits in a single data frame, which we can use for our statistical analysis.

\subsection{Analyzing and visualizing the results}

We can first visualize the relationship between semantic and biological similarity:

\begin{lstlisting}[style=Python]
sns.scatterplot(
    data=text_similarity_df,
    y='pathway_similarity',
    x='text_similarity',
    alpha=0.5,
    size=1
)
corr_value = text_similarity_df["pathway_similarity"].corr(
    text_similarity_df["text_similarity"])
plt.title(
  f'Pathway Similarity vs Text Similarity (r={corr_value:.2f})'
)
\end{lstlisting}

\begin{figure}[!htbp]
\centering
\includegraphics[width=0.625\linewidth]{files/pathway_vs_text_simi-c18ca6c7ca7a9d1e678437a11dd728ec.png}
\caption[Plot of biological versus semantic similarity]{A scatterplot of biological similarity (estimated as overlap in pathways) versus semantic similarity (estimated as embedding distance of PubMed abstracts) on the GWAS dataset.}
\label{PathwaySimilarity-fig}
\end{figure}

The result is show in Figure~\ref{PathwaySimilarity-fig}. There is a small but robust correlation between these two similarity measures. In order to more accurately estimate this association we need to take into account the fact that different phenotypes vary in their overall similarity by including a \textit{random effect} of phenotype within a mixed effects model.  We use the \texttt{lmer()} function from the R \texttt{lme4} package, via the R magic within Jupyter:\index{Jupyter}

\begin{lstlisting}[style=R]
%%R -i text_similarity_df

if (!requireNamespace("lme4", quietly = TRUE)) {
    install.packages("lme4")
}
if (!requireNamespace("lmerTest", quietly = TRUE)) {
    install.packages("lmerTest")
}
library(lme4)
library(lmerTest)

model < - lmer(pathway_similarity ~ text_similarity  + (1 | phenotype1) + (1 | phenotype2), data = text_similarity_df)
summary(model)
\end{lstlisting}

\begin{lstlisting}[style=shell]
Linear mixed model fit by REML. t -tests use Satterthwaite's method [
lmerModLmerTest]
Formula: pathway_similarity ~ text_similarity + (1 | phenotype1) + (1 |  
    phenotype2)
   Data: text_similarity_df

REML criterion at convergence: -9301.4

Scaled residuals: 
    Min      1Q  Median      3Q     Max 
 -5.4960 -0.4077 -0.0685  0.2873  8.9175 

Random effects:
 Groups     Name        Variance Std.Dev.
 phenotype1 (Intercept) 0.040351 0.20087 
 phenotype2 (Intercept) 0.009485 0.09739 
 Residual               0.007039 0.08390 
Number of obs: 6904, groups:  phenotype1, 981; phenotype2, 905

Fixed effects:
                 Estimate Std. Error        df t value Pr(>|t|)    
(Intercept)     2.484e -01  7.802e -03 1.507e+03   31.83   <2e -16 ***
text_similarity 3.666e -01  2.532e -02 5.858e+03   14.48   <2e -16 ***
 - - -
Signif. codes:  0 ‘***’ 0.001 ‘**’ 0.01 ‘*’ 0.05 ‘.’ 0.1 ‘ ’ 1

Correlation of Fixed Effects:
            (Intr)
text_smlrty -0.334
\end{lstlisting}

This result is statistically significant, but we also want to ask how practically significant it is by looking at the amount of variance accounted for by the model:

\begin{lstlisting}[style=R]
%%R

if (!requireNamespace("MuMIn", quietly = TRUE)) {
    install.packages("MuMIn")
}

library(MuMIn)
r.squaredGLMM(model)
\end{lstlisting}

\begin{lstlisting}
R2m       R2c
[1,] 0.01068788 0.8775626
\end{lstlisting}

For a mixed effect model there are two R-squared values: The \textit{conditional} R-squared (R2c), which refers to the total variance accounted for by both fixed and random effects, and the \textit{marginal} R-squared (R2m) that refers to the variance accounted for by the fixed effects, which is the measure of interest here.  This shows that while the association is strongly statistically significant, semantic similarity only accounts for about 1\% of the variability in biological similarity, and thus is not a particularly strong predictor in practice.  The high conditional R-squared demonstrates that the variability in the data is dominated by trait-level differences in similarity, such that traits vary in the degree to which they are more generally similar to others on average.